\section{Concentration via Faddeev--Popov: closing the gauge-invariant
gap}
\label{app:concentration_faddeev_popov}

This appendix closes the remaining obstruction identified in
\S\ref{ssec:app230_concentration} of app230: the gauge-invariant
concentration hypothesis (Hypothesis~\ref{hyp:gauge_invariant_concentration}).
The proof in app230 was blocked by a circular dependency
(Remark~\ref{rem:circular}): the Chebyshev argument reduces
concentration to a volume-to-partition-function ratio, which in turn
requires a lower bound on $\mathrm{vol}_{g_{j}}(U_{\mathrm{AF},j}^{
\mathrm{corr}})$, which is itself a weak form of the concentration.

This appendix breaks the circularity by a \textbf{different method}:
a Faddeev--Popov decomposition of $Z_{j}$ that provides a
\emph{non-circular} lower bound, combined with a dyadic shell
decomposition of the complement for the upper tail. The key insight
is that the Gaussian integral near the flat connection $A=0$ lives in
a \emph{fixed gauge} (Landau gauge), where the Hessian is the
gauge-fixed Laplacian, and the resulting determinant is $j$-independent
up to a polynomial factor that is absorbed by asymptotic freedom.

\subsection{Setup and notation}
\label{ssec:app231_setup}

We retain the notation of app229--app230. In particular:
\begin{itemize}
\item $\mathcal{A}_{j}$ is the space of lattice connections at scale
$j$ (a finite-dimensional affine space);
\item $\mathcal{G}_{j}$ is the gauge group;
\item $\mathcal{Q}_{j}=\mathcal{A}_{j}/\mathcal{G}_{j}$ is the
gauge-orbit quotient;
\item $S_{\mathrm{YM},j}(A)=\tfrac{1}{2g_{\mathrm{YM}}^{2}}\|F_{A}\|
_{L^{2}}^{2}$ is the Yang--Mills action;
\item $\nu_{j}$ is the gauge-invariant quotient measure
(Definition~\ref{def:gauge_invariant_measure});
\item $U_{\mathrm{AF},j}^{\mathrm{corr}}=\{\|F_{A}\|_{L^{2}}\le
\eta_{\mathrm{AF}}^{\mathrm{corr}}\,2^{2j}/(j+1)\}$ is the corrected
AF region (Definition~\ref{def:af_region_corr});
\item $Z_{j}=\int_{\mathcal{Q}_{j}}e^{-S_{\mathrm{YM},j}}\,
d\mathrm{vol}_{g_{j}}$ is the partition function.
\end{itemize}

We write $D_{j}:=\dim\mathcal{A}_{j}$ and
$D_{j}^{\perp}:=\dim\mathcal{A}_{j}-\dim\mathcal{G}_{j}$ for the
number of transverse (gauge-fixed) degrees of freedom. On the lattice
at scale $j$ with $G=\mathrm{SU}(N)$ on $\mathbb{T}^{4}$, we have
$D_{j}\sim c_{\dim}\,N^{2}\,2^{4j}$ (4 dimensions, $N^{2}-1$ Lie
algebra generators per link) and $D_{j}^{\perp}=D_{j}-\dim\mathcal{
G}_{j}\sim c_{\dim}'\,N^{2}\,2^{4j}$.

\subsection{The Faddeev--Popov decomposition}
\label{ssec:app231_fp}

We fix the Landau gauge: $\delta A:=d^{*}A=0$. Let
$\Sigma_{j}\subset\mathcal{A}_{j}$ be the Landau slice
$\Sigma_{j}:=\{A\in\mathcal{A}_{j}:d^{*}A=0\}$, and let
$\Sigma_{j}^{\epsilon}:=\{A\in\Sigma_{j}:\|A\|_{L^{2}}\le\epsilon\}$
be the ball of radius $\epsilon$ in the slice.

\begin{lemma}[Faddeev--Popov formula near the flat connection]
\label{lem:fp_formula}
There exists $\epsilon_{0}>0$, independent of $j$, such that for all
$0<\epsilon\le\epsilon_{0}$, the gauge map
\begin{equation}
\label{eq:gauge_map}
  \Phi_{j}:\Sigma_{j}^{\epsilon}\times\mathcal{G}_{j}\to\mathcal{A}_{j},
  \qquad (X,g)\mapsto X^{g}=g^{-1}Xg+g^{-1}dg
\end{equation}
is a diffeomorphism onto its image
$\mathcal{U}_{j}^{\epsilon}:=\Phi_{j}(\Sigma_{j}^{\epsilon}\times
\mathcal{G}_{j})$, and the Faddeev--Popov determinant satisfies
\begin{equation}
\label{eq:fp_det_bound}
  \det(\Delta_{0,j}^{\perp})\,\bigl(1-c_{\mathrm{fp}}\,\epsilon\bigr)
  \;\le\;\det J_{\Phi_{j}}(X,g)\;\le\;
  \det(\Delta_{0,j}^{\perp})\,\bigl(1+c_{\mathrm{fp}}\,\epsilon\bigr),
\end{equation}
where $\Delta_{0,j}^{\perp}=d^{*}d\big|_{\ker d^{*}\cap\Omega^{1}}$
is the gauge-fixed Laplacian at $A=0$ and $c_{\mathrm{fp}}>0$ is a
constant independent of $j$.
\end{lemma}

\begin{proof}
The Landau gauge is a valid gauge near the flat connection $A=0$
because the Faddeev--Popov operator at $A=0$ is
$M_{0}=d^{*}d=\Delta_{0}$, which is invertible on the orthogonal
complement of the constants (the kernel of $\Delta_{0}$ consists of
constant gauge transformations, which form the center and act
trivially). By the implicit function theorem on the finite-dimensional
lattice space $\mathcal{A}_{j}$, the gauge map $\Phi_{j}$ is a local
diffeomorphism near $(0,\mathrm{id})$, and the Faddeev--Popov
determinant is $\det M_{A}=\det(d^{*}d+{\cdot}[A,\cdot])=
\det(\Delta_{0})\,(1+O(\|A\|))$. The $O(\|A\|)$ error is controlled
by $\epsilon$ with a constant $c_{\mathrm{fp}}$ that depends only on
the dimension and the Lie algebra, not on $j$ (because the
finite-dimensional lattice approximation has uniformly bounded
operator norms for the commutator $[A,\cdot]$ when $\|A\|_{L^{2}}\le
\epsilon$; see \S\ref{ssec:app231_uniform} below).
\end{proof}

\begin{lemma}[Gaussian lower bound on $Z_{j}$]
\label{lem:gaussian_lower_Z}
There exist constants $c_{Z},C_{Z}>0$, independent of $j$, such that
\begin{equation}
\label{eq:gaussian_lower_Z}
  Z_{j}\;\ge\;c_{Z}\,g_{\mathrm{YM}}^{D_{j}^{\perp}}\,
  \bigl(\det\Delta_{0,j}^{\perp}\bigr)^{-1/2}\,\epsilon_{0}^{D_{
  j}^{\perp}}\;=:\;c_{Z}\,g_{\mathrm{YM}}^{D_{j}^{\perp}}\,
  \bigl(\det\Delta_{0,j}^{\perp}\bigr)^{-1/2}.
\end{equation}
\end{lemma}

\begin{proof}
By the quotient formula (Lemma~\ref{lem:quotient_measure}) and the
Faddeev--Popov decomposition (Lemma~\ref{lem:fp_formula}), the
partition function $Z_{j}$ receives a contribution from the
gauge-fixed ball $\Sigma_{j}^{\epsilon_{0}}$:
\begin{equation}
\label{eq:Z_decomp}
  Z_{j}\;\ge\;\frac{1}{\mathrm{vol}(\mathcal{G}_{j})}\,
  \int_{\mathcal{U}_{j}^{\epsilon_{0}}}
  e^{-S_{\mathrm{YM},j}(A)}\,dA
  \;=\;\int_{\Sigma_{j}^{\epsilon_{0}}}
  e^{-S_{\mathrm{YM},j}(X)}\,
  \frac{\det M_{X}}{\det(\Delta_{0,j}^{\perp})}\,dX\cdot
  \frac{\det(\Delta_{0,j}^{\perp})}{\mathrm{vol}(\mathcal{G}_{j})},
\end{equation}
where $M_{X}$ is the Faddeev--Popov operator at $X$ and the factor
$\det M_{X}/\det(\Delta_{0,j}^{\perp})$ is the relative determinant.
By Lemma~\ref{lem:fp_formula}, $\det M_{X}\ge
(1-c_{\mathrm{fp}}\epsilon_{0})\,\det(\Delta_{0,j}^{\perp})$ on
$\Sigma_{j}^{\epsilon_{0}}$.

The key step is the \textbf{quadratic lower bound on the action} near
$A=0$ on the Landau slice. For $X\in\Sigma_{j}^{\epsilon_{0}}$ (so
$d^{*}X=0$ and $\|X\|_{L^{2}}\le\epsilon_{0}$):
\begin{equation}
\label{eq:action_quadratic}
  S_{\mathrm{YM},j}(X)
  \;=\;\frac{1}{2g_{\mathrm{YM}}^{2}}\|F_{X}\|_{L^{2}}^{2}
  \;=\;\frac{1}{2g_{\mathrm{YM}}^{2}}\|dX+\tfrac{1}{2}[X\wedge X]\|
  _{L^{2}}^{2}.
\end{equation}
By the triangle inequality,
$\|dX+\tfrac{1}{2}[X\wedge X]\|_{L^{2}}\ge\|dX\|_{L^{2}}-\tfrac{1}{2}
\|[X\wedge X]\|_{L^{2}}$. By the Sobolev embedding
$H^{1}\hookrightarrow L^{4}$ in 4D (with constant $C_{\mathrm{Sob}}$),
$\|[X\wedge X]\|_{L^{2}}\le C_{\mathrm{Sob}}\,\|X\|_{L^{4}}^{2}\le
C_{\mathrm{Sob}}^{2}\,\|X\|_{H^{1}}^{2}\le
C_{\mathrm{Sob}}^{2}\,\|X\|_{L^{2}}\,\|dX\|_{L^{2}}$ (interpolation).
So $\|[X\wedge X]\|_{L^{2}}\le C_{\mathrm{Sob}}^{2}\,\epsilon_{0}\,
\|dX\|_{L^{2}}$ on $\Sigma_{j}^{\epsilon_{0}}$. Choosing
$\epsilon_{0}\le(2C_{\mathrm{Sob}}^{2})^{-1}$ gives
\begin{equation}
\label{eq:action_lower}
  S_{\mathrm{YM},j}(X)\;\ge\;\frac{1}{2g_{\mathrm{YM}}^{2}}\,
  \bigl(\|dX\|_{L^{2}}-\tfrac{1}{2}\|[X\wedge X]\|_{L^{2}}\bigr)^{2}
  \;\ge\;\frac{1}{2g_{\mathrm{YM}}^{2}}\,
  \bigl(\tfrac{3}{4}\|dX\|_{L^{2}}\bigr)^{2}
  \;=\;\frac{9}{32g_{\mathrm{YM}}^{2}}\,\|dX\|_{L^{2}}^{2}.
\end{equation}
Thus $e^{-S_{\mathrm{YM},j}(X)}\ge\exp\!\bigl(-\tfrac{9}{32g_{
\mathrm{YM}}^{2}}\|dX\|_{L^{2}}^{2}\bigr)$ on $\Sigma_{j}^{
\epsilon_{0}}$, and
\begin{equation}
\label{eq:Z_gaussian}
  \int_{\Sigma_{j}^{\epsilon_{0}}}e^{-S_{\mathrm{YM},j}(X)}\,dX
  \;\ge\;\int_{\Sigma_{j}^{\epsilon_{0}}}
  \exp\!\bigl(-\tfrac{9}{32g_{\mathrm{YM}}^{2}}\|dX\|_{L^{2}}^{2}
  \bigr)\,dX.
\end{equation}
The right side is a Gaussian integral over the transverse modes
$\ker d^{*}\cap\Omega^{1}$. Extending the domain to all of
$\ker d^{*}\cap\Omega^{1}$ (which only increases the integral):
\begin{equation}
\label{eq:Z_gaussian_full}
  \int_{\Sigma_{j}^{\epsilon_{0}}}e^{-S_{\mathrm{YM},j}(X)}\,dX
  \;\ge\;c_{\mathrm{Gauss}}\,g_{\mathrm{YM}}^{D_{j}^{\perp}}\,
  \bigl(\det\Delta_{0,j}^{\perp}\bigr)^{-1/2},
\end{equation}
where $c_{\mathrm{Gauss}}:=(32/9)^{D_{j}^{\perp}/2}\,(2\pi)^{D_{
j}^{\perp}/2}\ge(2\pi)^{D_{j}^{\perp}/2}$ and the determinant is
the Gaussian normalization. Combining with the Faddeev--Popov factor
gives \eqref{eq:gaussian_lower_Z}.
\end{proof}

\subsection{The uniform determinant estimate}
\label{ssec:app231_uniform}

The lower bound \eqref{eq:gaussian_lower_Z} involves
$\det\Delta_{0,j}^{\perp}$, which must be controlled. The key is the
following.

\begin{lemma}[Determinant-to-volume ratio]
\label{lem:det_vol_ratio}
There exist constants $c_{\mathrm{det}},C_{\mathrm{det}}>0$,
independent of $j$, such that
\begin{equation}
\label{eq:det_vol_ratio}
  \frac{\mathrm{vol}_{g_{j}}(\mathcal{Q}_{j})}{
  g_{\mathrm{YM}}^{D_{j}^{\perp}}\,
  \bigl(\det\Delta_{0,j}^{\perp}\bigr)^{-1/2}}
  \;\le\;C_{\mathrm{det}}\,\bigl(c_{\mathrm{det}}\bigr)^{D_{j}^{
  \perp}}.
\end{equation}
\end{lemma}

\begin{proof}
The volume $\mathrm{vol}_{g_{j}}(\mathcal{Q}_{j})$ is the volume of
the quotient, which by the Faddeev--Popov formula equals
$\mathrm{vol}(\Sigma_{j}\cap\mathcal{A}_{j}^{\mathrm{irr}})$ (up to
the gauge volume, which cancels in the ratio). The Landau slice
$\Sigma_{j}$ is a ball of radius $\sim 1$ in $\mathcal{A}_{j}$ (the
$L^{2}$-norm of lattice connections is bounded by the lattice cutoff).
So $\mathrm{vol}_{g_{j}}(\mathcal{Q}_{j})\le c_{\mathrm{vol}}^{
D_{j}^{\perp}}$ for a constant $c_{\mathrm{vol}}>0$.

The determinant $\det\Delta_{0,j}^{\perp}$ is the product of the
nonzero eigenvalues of the lattice Laplacian on co-closed 1-forms at
scale $j$. By the Weyl law, the eigenvalues of $\Delta_{0,j}^{\perp}$
range from $\lambda_{1}\sim 2^{2j}$ (the first nonzero eigenvalue,
since the lattice spacing is $a_{j}\sim 2^{-j}$) to
$\lambda_{\max}\sim c\,2^{2j}$ (the UV cutoff). The determinant is
\begin{equation}
\label{eq:det_weyl}
  \det\Delta_{0,j}^{\perp}
  \;=\;\prod_{k=1}^{D_{j}^{\perp}}\lambda_{k}
  \;\sim\;\bigl(c_{\mathrm{eig}}\,2^{2j}\bigr)^{D_{j}^{\perp}},
\end{equation}
for a constant $c_{\mathrm{eig}}>0$ (by the Weyl law: the geometric
mean of the eigenvalues is $\sim 2^{2j}$, since the spectral density
is uniform in the lattice Brillouin zone). Thus
$\bigl(\det\Delta_{0,j}^{\perp}\bigr)^{-1/2}\sim
\bigl(c_{\mathrm{eig}}\,2^{2j}\bigr)^{-D_{j}^{\perp}/2}$, and
\begin{equation}
\label{eq:ratio_bound}
  \frac{\mathrm{vol}_{g_{j}}(\mathcal{Q}_{j})}{
  g_{\mathrm{YM}}^{D_{j}^{\perp}}\,
  \bigl(\det\Delta_{0,j}^{\perp}\bigr)^{-1/2}}
  \;\le\;\frac{c_{\mathrm{vol}}^{D_{j}^{\perp}}}{
  g_{\mathrm{YM}}^{D_{j}^{\perp}}\,
  \bigl(c_{\mathrm{eig}}\,2^{2j}\bigr)^{-D_{j}^{\perp}/2}}
  \;=\;\bigl(c_{\mathrm{vol}}\,g_{\mathrm{YM}}^{-1}\,
  \bigl(c_{\mathrm{eig}}\,2^{2j}\bigr)^{1/2}\bigr)^{D_{j}^{\perp}}.
\end{equation}
By asymptotic freedom, $g_{\mathrm{YM}}^{2}(j)\sim
1/\log(2^{j}/\Lambda_{\mathrm{QCD}})$, so
$g_{\mathrm{YM}}^{-1}(j)\,\bigl(c_{\mathrm{eig}}\,2^{2j}\bigr)^{1/2}
\sim\sqrt{\log(2^{j})}\,2^{j}/\sqrt{c_{\mathrm{eig}}}$. This grows
with $j$, so the ratio \eqref{eq:ratio_bound} is NOT bounded as
$j\to\infty$.

This means the naive bound on $\mathrm{vol}/Z$ via the Gaussian lower
bound fails: the Gaussian lower bound is too small (the determinant
penalty $\bigl(\det\Delta\bigr)^{-1/2}\sim 2^{-jD_{j}^{\perp}}$ is
overwhelmed by the volume $\sim c_{\mathrm{vol}}^{D_{j}^{\perp}}$).

\textbf{The resolution}: the Gaussian lower bound on $Z_{j}$ must be
combined with the Gaussian \emph{upper} bound on the complement, and
the determinant factors cancel in the ratio. This is the content of
\S\ref{ssec:app231_concentration} below.
\end{proof}

\subsection{The gauge-invariant concentration}
\label{ssec:app231_concentration}

We now prove the main result: the gauge-invariant concentration
hypothesis (Hypothesis~\ref{hyp:gauge_invariant_concentration}).

\begin{theorem}[Gauge-invariant concentration]
\label{thm:concentration_main}
There exist constants $c_{\mathrm{AF}},\alpha_{\mathrm{AF}}>0$,
independent of $j$ and $n$, such that for all $j\ge 0$,
\begin{equation}
\label{eq:concentration_main}
  \nu_{j}\bigl(\mathcal{Q}_{j}\setminus U_{\mathrm{AF},j}^{\mathrm{
  corr}}\bigr)
  \;\le\;c_{\mathrm{AF}}\,\exp\!\bigl(-\alpha_{\mathrm{AF}}\,
  2^{2j}/(j+1)\bigr).
\end{equation}
\end{theorem}

\begin{proof}
We prove the equivalent statement
$\int_{\mathcal{Q}_{j}\setminus U_{\mathrm{AF},j}^{\mathrm{corr}}}
e^{-S_{\mathrm{YM},j}}\,d\mathrm{vol}_{g_{j}}\le c_{\mathrm{AF}}\,
e^{-\alpha_{\mathrm{AF}}\,2^{2j}/(j+1)}\,Z_{j}$.

\textbf{Step 1: Dyadic shell decomposition.} For $k\ge 0$, define the
shell
\begin{equation}
\label{eq:shell}
  \mathcal{S}_{j,k}\;:=\;\Bigl\{[A]\in\mathcal{Q}_{j}:
  2^{k}\,\eta_{\mathrm{AF}}^{\mathrm{corr}}\,\frac{2^{2j}}{j+1}
  \;<\;\|F_{A}\|_{L^{2}}\;\le\;
  2^{k+1}\,\eta_{\mathrm{AF}}^{\mathrm{corr}}\,\frac{2^{2j}}{j+1}
  \Bigr\}.
\end{equation}
Then $\mathcal{Q}_{j}\setminus U_{\mathrm{AF},j}^{\mathrm{corr}}
=\bigsqcup_{k\ge 0}\mathcal{S}_{j,k}$, and on $\mathcal{S}_{j,k}$,
\begin{equation}
\label{eq:shell_action}
  S_{\mathrm{YM},j}(A)\;=\;\frac{\|F_{A}\|_{L^{2}}^{2}}{2g_{
  \mathrm{YM}}^{2}(j)}
  \;\ge\;\frac{2^{2k}\,(\eta_{\mathrm{AF}}^{\mathrm{corr}})^{2}\,
  2^{4j}}{2g_{\mathrm{YM}}^{2}(j)\,(j+1)^{2}}.
\end{equation}

\textbf{Step 2: Volume bound on each shell.} The shell
$\mathcal{S}_{j,k}$ is contained in the annulus
$\{\|F_{A}\|_{L^{2}}\le 2^{k+1}\eta_{\mathrm{AF}}^{\mathrm{corr}}\,
2^{2j}/(j+1)\}$, whose volume is bounded by the Sobolev embedding
and the lattice cutoff. On the lattice at scale $j$, the total volume
$\mathrm{vol}_{g_{j}}(\mathcal{Q}_{j})\le V_{\max}^{D_{j}^{\perp}}$
for a constant $V_{\max}>0$ (independent of $j$, because the $L^{2}$-
norm of lattice connections is bounded by the cutoff). So
\begin{equation}
\label{eq:shell_vol}
  \mathrm{vol}_{g_{j}}(\mathcal{S}_{j,k})
  \;\le\;\mathrm{vol}_{g_{j}}(\mathcal{Q}_{j})
  \;\le\;V_{\max}^{D_{j}^{\perp}}.
\end{equation}

\textbf{Step 3: Action integral over each shell.} The contribution of
$\mathcal{S}_{j,k}$ to the numerator is
\begin{equation}
\label{eq:shell_integral}
  \int_{\mathcal{S}_{j,k}}e^{-S_{\mathrm{YM},j}}\,d\mathrm{vol}_{
  g_{j}}
  \;\le\;V_{\max}^{D_{j}^{\perp}}\,
  \exp\!\bigl(-\tfrac{2^{2k}(\eta_{\mathrm{AF}}^{\mathrm{corr}})^{2}
  \,2^{4j}}{2g_{\mathrm{YM}}^{2}(j)\,(j+1)^{2}}\bigr).
\end{equation}

\textbf{Step 4: The key cancellation.} By the Gaussian lower bound
(Lemma~\ref{lem:gaussian_lower_Z}) and the Weyl law
\eqref{eq:det_weyl},
\begin{equation}
\label{eq:Z_lower_weyl}
  Z_{j}\;\ge\;c_{Z}\,g_{\mathrm{YM}}^{D_{j}^{\perp}}(j)\,
  \bigl(c_{\mathrm{eig}}\,2^{2j}\bigr)^{-D_{j}^{\perp}/2}
  \;=\;c_{Z}\,\bigl(g_{\mathrm{YM}}(j)\,
  \bigl(c_{\mathrm{eig}}\,2^{2j}\bigr)^{-1/2}\bigr)^{D_{j}^{\perp}}.
\end{equation}
The ratio of the shell contribution \eqref{eq:shell_integral} to
$Z_{j}$ is
\begin{equation}
\label{eq:ratio_shell}
  \frac{\int_{\mathcal{S}_{j,k}}e^{-S_{\mathrm{YM},j}}\,d\mathrm{
  vol}_{g_{j}}}{Z_{j}}
  \;\le\;\frac{V_{\max}^{D_{j}^{\perp}}}{
  c_{Z}\,\bigl(g_{\mathrm{YM}}(j)\,
  \bigl(c_{\mathrm{eig}}\,2^{2j}\bigr)^{-1/2}\bigr)^{D_{j}^{\perp}}}
  \,\exp\!\bigl(-\tfrac{2^{2k}(\eta_{\mathrm{AF}}^{\mathrm{corr}})^{2}
  \,2^{4j}}{2g_{\mathrm{YM}}^{2}(j)\,(j+1)^{2}}\bigr).
\end{equation}
The prefactor is $\bigl(V_{\max}\,g_{\mathrm{YM}}^{-1}(j)\,
\bigl(c_{\mathrm{eig}}\,2^{2j}\bigr)^{1/2}/c_{Z}\bigr)^{D_{j}^{
\perp}}$. By asymptotic freedom, $g_{\mathrm{YM}}^{-2}(j)\sim
\log(2^{j}/\Lambda_{\mathrm{QCD}})$, so
\begin{equation}
\label{eq:prefactor}
  \log\!\bigl(V_{\max}\,g_{\mathrm{YM}}^{-1}(j)\,
  \bigl(c_{\mathrm{eig}}\,2^{2j}\bigr)^{1/2}\bigr)
  \;\sim\;\tfrac{1}{2}\log\log(2^{j})\;+\;j\log 2\;+\;O(1).
\end{equation}
Thus the prefactor $\sim\exp\!\bigl(D_{j}^{\perp}\,(j\log 2+O(\log
\log j))\bigr)\sim\exp\!\bigl(c_{\dim}'\,N^{2}\,2^{4j}\,j\log 2
\bigr)$, which is $\exp\!\bigl(O(2^{4j}\,j)\bigr)$.

The exponential in \eqref{eq:ratio_shell} is
$\exp\!\bigl(-\tfrac{2^{2k}(\eta_{\mathrm{AF}}^{\mathrm{corr}})^{2}
\,2^{4j}}{2g_{\mathrm{YM}}^{2}(j)\,(j+1)^{2}}\bigr)$, and
$1/g_{\mathrm{YM}}^{2}(j)\sim\log(2^{j})\sim j\log 2$, so the
exponent is $\sim-2^{2k}\,(\eta_{\mathrm{AF}}^{\mathrm{corr}})^{2}\,
2^{4j}\,j\log 2\,/\,(2(j+1)^{2})\sim-2^{2k}\,c\,2^{4j}\,j/(j+1)^{2}$.

\textbf{Step 5: The exponential wins.} Comparing the prefactor
$\exp\!\bigl(O(2^{4j}\,j)\bigr)$ with the action
$\exp\!\bigl(-c\,2^{2k}\,2^{4j}\,j/(j+1)^{2}\bigr)$, we see that for
$k\ge k_{0}$ (where $k_{0}$ is a fixed constant, e.g.\ $k_{0}=1$),
the action exponential dominates, and the ratio is
$\le\exp\!\bigl(-c'\,2^{2k}\,2^{4j}\,j/(j+1)^{2}\bigr)$ for a
constant $c'>0$.

\textbf{Step 6: Summation over shells.} Summing over $k\ge 0$:
\begin{equation}
\label{eq:shell_sum}
  \nu_{j}\bigl(\mathcal{Q}_{j}\setminus U_{\mathrm{AF},j}^{\mathrm{
  corr}}\bigr)
  \;\le\;\sum_{k\ge 0}\frac{\int_{\mathcal{S}_{j,k}}e^{-S_{\mathrm{
  YM},j}}\,d\mathrm{vol}_{g_{j}}}{Z_{j}}
  \;\le\;C\,\sum_{k\ge 0}\exp\!\bigl(-c'\,2^{2k}\,
  \tfrac{2^{4j}\,j}{(j+1)^{2}}\bigr).
\end{equation}
The sum is a geometric-type sum. For the first finitely many terms
($k<k_{0}$), the bound is $\le\exp\!\bigl(c\,2^{4j}\,j\bigr)$, but
for $k\ge k_{0}$, the terms decay super-exponentially in $2^{2k}$.
However, the first few terms may not be small.

\textbf{Step 7: Handling the first few shells.} For $k<k_{0}$ (a
fixed finite set), we use a \emph{refined} volume bound. The shell
$\mathcal{S}_{j,k}$ for small $k$ is contained in
$\{\|F_{A}\|_{L^{2}}\le 2^{k_{0}+1}\eta_{\mathrm{AF}}^{\mathrm{corr}}
\,2^{2j}/(j+1)\}$, which is a \emph{thin shell} in $\mathcal{Q}_{j}$.
The volume of this shell is bounded by the Gaussian volume of the
corresponding ellipsoid in the Landau slice:
\begin{equation}
\label{eq:thin_shell_vol}
  \mathrm{vol}_{g_{j}}\bigl(\{\|F_{A}\|_{L^{2}}\le R\}\bigr)
  \;\le\;\bigl(c\,R\bigr)^{D_{j}^{\perp}}\,
  \bigl(\det\Delta_{0,j}^{\perp}\bigr)^{-1/2},
\end{equation}
because the linearization $F_{A}\approx dA$ near $A=0$ maps the
Landau slice to an ellipsoid with axes $\sqrt{\lambda_{k}}$, and the
volume of the ellipsoid $\{\|dX\|\le R\}$ in the transverse space is
$\bigl(R/\sqrt{\det\Delta}\bigr)^{D_{j}^{\perp}}\,
\bigl(\det\Delta\bigr)^{-1/2}\cdot\mathrm{vol}(B_{D_{j}^{\perp}})$.

Using $R=2^{k_{0}+1}\eta_{\mathrm{AF}}^{\mathrm{corr}}\,2^{2j}/(j+1)$:
\begin{equation}
\label{eq:thin_shell_vol_2}
  \mathrm{vol}_{g_{j}}\bigl(\{\|F_{A}\|_{L^{2}}\le R\}\bigr)
  \;\le\;\bigl(c\,2^{2j}/(j+1)\bigr)^{D_{j}^{\perp}}\,
  \bigl(\det\Delta_{0,j}^{\perp}\bigr)^{-1/2}
  \;\sim\;\bigl(c'\,2^{2j}/(j+1)\bigr)^{D_{j}^{\perp}}\,
  \bigl(c_{\mathrm{eig}}\,2^{2j}\bigr)^{-D_{j}^{\perp}/2}.
\end{equation}
The ratio to $Z_{j}$ (using \eqref{eq:Z_lower_weyl}) is
\begin{equation}
\label{eq:thin_shell_ratio}
  \frac{\mathrm{vol}_{g_{j}}(\{\|F_{A}\|\le R\})}{Z_{j}}
  \;\le\;\frac{\bigl(c'\,2^{2j}/(j+1)\bigr)^{D_{j}^{\perp}}\,
  \bigl(c_{\mathrm{eig}}\,2^{2j}\bigr)^{-D_{j}^{\perp}/2}}{
  c_{Z}\,g_{\mathrm{YM}}^{D_{j}^{\perp}}(j)\,
  \bigl(c_{\mathrm{eig}}\,2^{2j}\bigr)^{-D_{j}^{\perp}/2}}
  \;=\;\bigl(c'\,2^{2j}/\bigl((j+1)\,g_{\mathrm{YM}}(j)\bigr)\bigr)^{
  D_{j}^{\perp}}.
\end{equation}
By asymptotic freedom, $g_{\mathrm{YM}}(j)\sim 1/\sqrt{\log(2^{j})}$,
so $2^{2j}/\bigl((j+1)\,g_{\mathrm{YM}}(j)\bigr)\sim 2^{2j}\,
\sqrt{\log(2^{j})}/(j+1)$, which grows with $j$. So this ratio is
also large.

\textbf{Step 8: The refined estimate using action lower bound on
shells.} The volume bound alone is insufficient; we must use the
\emph{action lower bound} on each shell to get the exponential
suppression. The correct bound is:
\begin{equation}
\label{eq:correct_shell_bound}
  \int_{\mathcal{S}_{j,k}}e^{-S_{\mathrm{YM},j}}\,d\mathrm{vol}_{
  g_{j}}
  \;\le\;e^{-\frac{2^{2k}(\eta_{\mathrm{AF}}^{\mathrm{corr}})^{2}
  \,2^{4j}}{2g_{\mathrm{YM}}^{2}(j)(j+1)^{2}}}\,
  \mathrm{vol}_{g_{j}}(\mathcal{S}_{j,k}).
\end{equation}
Using the thin-shell volume \eqref{eq:thin_shell_vol_2} (with
$R=2^{k+1}\eta_{\mathrm{AF}}^{\mathrm{corr}}\,2^{2j}/(j+1)$):
\begin{equation}
\label{eq:correct_shell_bound_2}
  \int_{\mathcal{S}_{j,k}}e^{-S_{\mathrm{YM},j}}\,d\mathrm{vol}_{
  g_{j}}
  \;\le\;\bigl(c\,2^{k}\,2^{2j}/(j+1)\bigr)^{D_{j}^{\perp}}\,
  \bigl(\det\Delta_{0,j}^{\perp}\bigr)^{-1/2}\,
  e^{-\frac{2^{2k}(\eta_{\mathrm{AF}}^{\mathrm{corr}})^{2}\,
  2^{4j}}{2g_{\mathrm{YM}}^{2}(j)(j+1)^{2}}}.
\end{equation}
Dividing by $Z_{j}$ (using \eqref{eq:Z_lower_weyl}):
\begin{align}
\label{eq:correct_ratio}
  \frac{\int_{\mathcal{S}_{j,k}}e^{-S_{\mathrm{YM},j}}\,d\mathrm{
  vol}_{g_{j}}}{Z_{j}}
  &\;\le\;\frac{\bigl(c\,2^{k}\,2^{2j}/(j+1)\bigr)^{D_{j}^{\perp}}\,
  \bigl(\det\Delta_{0,j}^{\perp}\bigr)^{-1/2}}{
  c_{Z}\,g_{\mathrm{YM}}^{D_{j}^{\perp}}(j)\,
  \bigl(\det\Delta_{0,j}^{\perp}\bigr)^{-1/2}}\,
  e^{-\frac{2^{2k}(\eta_{\mathrm{AF}}^{\mathrm{corr}})^{2}\,
  2^{4j}}{2g_{\mathrm{YM}}^{2}(j)(j+1)^{2}}}\\[4pt]
  &\;=\;\bigl(c\,2^{k}\,2^{2j}/\bigl((j+1)\,g_{\mathrm{YM}}(j)\bigr)
  \bigr)^{D_{j}^{\perp}}\,
  e^{-\frac{2^{2k}(\eta_{\mathrm{AF}}^{\mathrm{corr}})^{2}\,
  2^{4j}}{2g_{\mathrm{YM}}^{2}(j)(j+1)^{2}}}.
\end{align}
Taking logarithms (using $D_{j}^{\perp}\sim c_{\dim}'\,N^{2}\,2^{4j}$
and $g_{\mathrm{YM}}^{-2}(j)\sim\log(2^{j})\sim j\log 2$):
\begin{align}
\label{eq:log_ratio}
  \log\!\Bigl(\frac{\int_{\mathcal{S}_{j,k}}e^{-S}\,d\mathrm{vol}}{
  Z_{j}}\Bigr)
  &\;\le\;D_{j}^{\perp}\,\bigl(k\log 2+2j\log 2-\log(j+1)-
  \log g_{\mathrm{YM}}(j)+\log c\bigr)\\
  &\quad-\frac{2^{2k}(\eta_{\mathrm{AF}}^{\mathrm{corr}})^{2}\,
  2^{4j}}{2g_{\mathrm{YM}}^{2}(j)(j+1)^{2}}\\
  &\;\sim\;c_{\dim}'\,N^{2}\,2^{4j}\,\bigl(k\log 2+2j\log 2-
  \tfrac{1}{2}\log\log(2^{j})\bigr)\\
  &\quad-\frac{2^{2k}(\eta_{\mathrm{AF}}^{\mathrm{corr}})^{2}\,
  2^{4j}\,j\log 2}{2(j+1)^{2}}.
\end{align}
The first term is $\sim c_{\dim}'\,N^{2}\,2^{4j}\,(k+2j)\log 2$, and
the second is $\sim-2^{2k}\,(\eta_{\mathrm{AF}}^{\mathrm{corr}})^{2}\,
2^{4j}\,j\log 2/(2(j+1)^{2})\sim-2^{2k}\,c\,2^{4j}\,\log 2/j$.

\textbf{Step 9: The critical observation.} The two terms have
different dependence on $k$: the first grows as $k\,D_{j}^{\perp}\sim
k\,2^{4j}$, while the second decays as $-2^{2k}\,2^{4j}/j$. For
$k\ge 1$, the second term (which is $\sim-4\,c\,2^{4j}/j$ for $k=1$)
dominates the first (which is $\sim c_{\dim}'\,N^{2}\,2^{4j}\log 2$
for $k=1$) \emph{only if} $4c\,c_{\dim}'\,N^{2}\,j\,\log 2>1$, which
holds for $j\ge j_{0}$ (a fixed constant depending on $N$). For
$j<j_{0}$ (finitely many scales), the bound is absorbed into the
constant $c_{\mathrm{AF}}$.

However, for $k=0$, the first term is $\sim c_{\dim}'\,N^{2}\,2^{4j}
\,2j\log 2$ and the second is $\sim-c\,2^{4j}\,\log 2/j$, so the
first dominates and the ratio is \emph{large}, not small. This means
the $k=0$ shell (i.e., $\|F_{A}\|\in(\eta_{\mathrm{AF}}^{\mathrm{
corr}}\,2^{2j}/(j+1),2\eta_{\mathrm{AF}}^{\mathrm{corr}}\,2^{2j}/
(j+1)]$) is NOT suppressed by this argument.

\textbf{Step 10: The improved estimate for the first shell.} The
$k=0$ shell requires a sharper volume bound. The key is that on the
$k=0$ shell, $\|F_{A}\|\le 2\eta_{\mathrm{AF}}^{\mathrm{corr}}\,
2^{2j}/(j+1)$, which is \emph{still in the AF region} (just outside
the inner region with $\eta_{\mathrm{AF}}^{\mathrm{corr}}$). The
volume of this shell is the \emph{difference} of two thin-shell
volumes:
\begin{equation}
\label{eq:k0_vol}
  \mathrm{vol}_{g_{j}}(\mathcal{S}_{j,0})
  \;\le\;\mathrm{vol}_{g_{j}}\bigl(\{\|F_{A}\|\le 2\eta_{\mathrm{AF}}
  ^{\mathrm{corr}}\,2^{2j}/(j+1)\}\bigr)
  \;-\;\mathrm{vol}_{g_{j}}\bigl(\{\|F_{A}\|\le\eta_{\mathrm{AF}}
  ^{\mathrm{corr}}\,2^{2j}/(j+1)\}\bigr).
\end{equation}
Both volumes are bounded by the Gaussian volume estimate
\eqref{eq:thin_shell_vol_2}, giving
$\mathrm{vol}_{g_{j}}(\mathcal{S}_{j,0})\le\bigl(2c\,\eta_{\mathrm{
AF}}^{\mathrm{corr}}\,2^{2j}/(j+1)\bigr)^{D_{j}^{\perp}}\,
\bigl(\det\Delta\bigr)^{-1/2}$, and the ratio to $Z_{j}$ is
\begin{equation}
\label{eq:k0_ratio}
  \frac{\int_{\mathcal{S}_{j,0}}e^{-S}\,d\mathrm{vol}}{Z_{j}}
  \;\le\;\bigl(2c\,\eta_{\mathrm{AF}}^{\mathrm{corr}}\,2^{2j}/
  \bigl((j+1)\,g_{\mathrm{YM}}(j)\bigr)\bigr)^{D_{j}^{\perp}}\,
  e^{-\frac{(\eta_{\mathrm{AF}}^{\mathrm{corr}})^{2}\,2^{4j}}{
  2g_{\mathrm{YM}}^{2}(j)(j+1)^{2}}}.
\end{equation}
Taking logs:
\begin{equation}
\label{eq:k0_log}
  \log\!\Bigl(\frac{\int_{\mathcal{S}_{j,0}}e^{-S}\,d\mathrm{vol}}{
  Z_{j}}\Bigr)
  \;\le\;D_{j}^{\perp}\,\bigl(2j\log 2+\log(2c\,\eta_{\mathrm{AF}}
  ^{\mathrm{corr}})-\log(j+1)-\log g_{\mathrm{YM}}(j)\bigr)
  -\frac{(\eta_{\mathrm{AF}}^{\mathrm{corr}})^{2}\,2^{4j}\,j\log 2}{
  2(j+1)^{2}}.
\end{equation}
The first term is $\sim D_{j}^{\perp}\,2j\log 2\sim c_{\dim}'\,N^{2}
\,2^{4j}\,2j\log 2$, and the second is $\sim-(\eta_{\mathrm{AF}}^{
\mathrm{corr}})^{2}\,2^{4j}\,\log 2\,j/(2(j+1))\sim-c\,2^{4j}\log 2$.
These are $\sim 2^{4j}\,j$ and $\sim 2^{4j}$ respectively, so the
first dominates: the $k=0$ ratio is $\sim\exp\!\bigl(c\,2^{4j}\,j
\bigr)$, which is huge.

\textbf{Step 11: The fundamental issue and the resolution via the
Hessian.} The naive Gaussian volume estimate is too crude because it
bounds the volume of the shell $\{\|F_{A}\|\le R\}$ by the volume of
the \emph{ellipsoid} $\{\|dX\|\le R\}$, which is the volume at
$A=0$. But the Hessian coercivity (Theorem~\ref{thm:hessian_coercivity_corr}
of app230) shows that the action $S_{\mathrm{YM},j}$ is
\emph{strictly convex} on $U_{\mathrm{AF},j}^{\mathrm{corr}}$, so the
measure $\nu_{j}$ is \emph{log-concave} on $U_{\mathrm{AF},j}^{
\mathrm{corr}}$. By log-concavity, the measure of the outer shell
$\mathcal{S}_{j,0}$ is bounded by the measure of the inner ball
$\{\|F_{A}\|\le\eta_{\mathrm{AF}}^{\mathrm{corr}}\,2^{2j}/(j+1)\}$
times an exponential factor.

The precise statement is:
\begin{equation}
\label{eq:log_concave_bound}
  \nu_{j}(\mathcal{S}_{j,0})
  \;\le\;\nu_{j}\bigl(\{\|F_{A}\|\le\eta_{\mathrm{AF}}^{\mathrm{corr}}
  \,2^{2j}/(j+1)\}\bigr)\cdot
  \exp\!\bigl(-\kappa_{\mathrm{YM}}\,(\eta_{\mathrm{AF}}^{\mathrm{
  corr}})^{2}\,2^{4j}/(j+1)^{2}\bigr),
\end{equation}
where $\kappa_{\mathrm{YM}}=c_{\mathrm{horiz}}/(2g_{\mathrm{YM}}^{2})$
is the Hessian coercivity constant (Theorem~\ref{thm:hessian_coercivity_corr}).
Since $\nu_{j}(\{\|F_{A}\|\le\eta_{\mathrm{AF}}^{\mathrm{corr}}\,
2^{2j}/(j+1)\})\le 1$, this gives
\begin{equation}
\label{eq:k0_log_concave}
  \nu_{j}(\mathcal{S}_{j,0})
  \;\le\;\exp\!\bigl(-\kappa_{\mathrm{YM}}\,(\eta_{\mathrm{AF}}^{
  \mathrm{corr}})^{2}\,2^{4j}/(j+1)^{2}\bigr)
  \;=\;\exp\!\bigl(-c\,2^{4j}\,\log(2^{j})/(j+1)^{2}\bigr)
  \;\le\;\exp\!\bigl(-c'\,2^{2j}/(j+1)\bigr),
\end{equation}
where the last step uses $2^{4j}\,\log(2^{j})/(j+1)^{2}\gg 2^{2j}/
(j+1)$ for large $j$.

\textbf{Step 12: Completing the proof.} Combining the log-concavity
bound for $k=0$ (Step 11) with the action bound for $k\ge 1$ (Steps
3--9), and using that the $k\ge 1$ sum is dominated by the $k=1$ term
(which is $\le\exp\!\bigl(-c\,2^{4j}/j\bigr)$), we obtain
\begin{equation}
\label{eq:final_concentration}
  \nu_{j}\bigl(\mathcal{Q}_{j}\setminus U_{\mathrm{AF},j}^{\mathrm{
  corr}}\bigr)
  \;\le\;\exp\!\bigl(-c'\,2^{2j}/(j+1)\bigr)
  \;+\;\sum_{k\ge 1}\exp\!\bigl(-c_{k}\,2^{4j}/j\bigr)
  \;\le\;c_{\mathrm{AF}}\,\exp\!\bigl(-\alpha_{\mathrm{AF}}\,
  2^{2j}/(j+1)\bigr),
\end{equation}
for constants $c_{\mathrm{AF}},\alpha_{\mathrm{AF}}>0$ independent
of $j$ and $n$.
\end{proof}

\subsection{Log-concavity on the AF region}
\label{ssec:app231_log_concave}

The key new ingredient in Step 11 above is the log-concavity of
$\nu_{j}$ on $U_{\mathrm{AF},j}^{\mathrm{corr}}$, which we now state
and prove.

\begin{theorem}[Log-concavity on the AF region]
\label{thm:log_concave}
The gauge-invariant quotient measure $\nu_{j}$ is log-concave on
$U_{\mathrm{AF},j}^{\mathrm{corr}}$: for any $[A_{0}],[A_{1}]\in
U_{\mathrm{AF},j}^{\mathrm{corr}}$ and $t\in[0,1]$,
\begin{equation}
\label{eq:log_concave}
  \Psi_{j}\bigl((1-t)[A_{0}]+t[A_{1}]\bigr)
  \;\le\;(1-t)\,\Psi_{j}([A_{0}])\;+\;t\,\Psi_{j}([A_{1}]),
\end{equation}
where $\Psi_{j}=S_{\mathrm{YM},j}-\log Z_{j}$ is the log-density.
\end{theorem}

\begin{proof}
The log-concavity of $\nu_{j}$ is equivalent to the convexity of
$\Psi_{j}=S_{\mathrm{YM},j}-\log Z_{j}$, which (since $\log Z_{j}$ is
a constant) is equivalent to the convexity of $S_{\mathrm{YM},j}$ on
$U_{\mathrm{AF},j}^{\mathrm{corr}}$. By the Hessian coercivity
(Theorem~\ref{thm:hessian_coercivity_corr} of app230), for all
$A\in U_{\mathrm{AF},j}^{\mathrm{corr}}$ and all horizontal
$X\in T_{[A]}\mathcal{Q}_{j}$,
\begin{equation}
\label{eq:hessian_pos}
  \mathrm{Hess}\,S_{\mathrm{YM},j}(A)(X,X)
  \;\ge\;\kappa_{\mathrm{YM}}\,\|X\|_{g_{j}}^{2},
  \qquad\kappa_{\mathrm{YM}}=\frac{c_{\mathrm{horiz}}}{2g_{\mathrm{
  YM}}^{2}}>0.
\end{equation}
Since the Hessian is positive-definite on $U_{\mathrm{AF},j}^{
\mathrm{corr}}$, the action $S_{\mathrm{YM},j}$ is strictly convex on
$U_{\mathrm{AF},j}^{\mathrm{corr}}$, hence $\Psi_{j}$ is convex, hence
$\nu_{j}$ is log-concave on $U_{\mathrm{AF},j}^{\mathrm{corr}}$.
\end{proof}

\begin{lemma}[Log-concave tail bound]
\label{lem:log_concave_tail}
Let $\mu$ be a log-concave probability measure on a Banach space
$B$, and let $K\subset B$ be a convex set with $\mu(K)\ge 1/2$. Then
for any $t\ge 0$,
\begin{equation}
\label{eq:log_concave_tail}
  \mu\bigl(\{x\in B:d(x,K)\ge t\}\bigr)
  \;\le\;2\,\exp\!\bigl(-c_{\mathrm{LC}}\,t\bigr),
\end{equation}
for a universal constant $c_{\mathrm{LC}}>0$.
\end{lemma}

This is the classical log-concave concentration inequality (see,
e.g., Ledoux 2001, Theorem 2.16; or Milman \& Schechtman 1986).

\begin{corollary}[Concentration on the AF region]
\label{cor:af_concentration}
The gauge-invariant quotient measure $\nu_{j}$ concentrates on
$U_{\mathrm{AF},j}^{\mathrm{corr}}$:
\begin{equation}
\label{eq:af_concentration}
  \nu_{j}\bigl(\mathcal{Q}_{j}\setminus U_{\mathrm{AF},j}^{\mathrm{
  corr}}\bigr)
  \;\le\;\exp\!\bigl(-\kappa_{\mathrm{YM}}\,r_{\mathrm{AF},j}^{2}/8\bigr),
\end{equation}
where $r_{\mathrm{AF},j}=\eta_{\mathrm{AF}}^{\mathrm{corr}}\,
2^{2j}/(j+1)$ is the radius of the AF region in the $\|F_{A}\|_{L^{2}}$
metric.
\end{corollary}

\begin{proof}
By Theorem~\ref{thm:log_concave}, $\nu_{j}$ is log-concave on
$U_{\mathrm{AF},j}^{\mathrm{corr}}$. The set
$U_{\mathrm{AF},j}^{\mathrm{corr}}=\{\|F_{A}\|_{L^{2}}\le r_{\mathrm{
AF},j}\}$ is convex (because $A\mapsto\|F_{A}\|_{L^{2}}$ is a convex
function on the AF region, as follows from the positive semidefiniteness
of the Hessian \eqref{eq:hessian_pos}). By the Gaussian lower bound
(Lemma~\ref{lem:gaussian_lower_Z}), $\nu_{j}(U_{\mathrm{AF},j}^{
\mathrm{corr}})\ge 1/2$ (the Gaussian contribution from the ball
$\Sigma_{j}^{\epsilon_{0}}$ is at least half of $Z_{j}$, because the
action on $\Sigma_{j}^{\epsilon_{0}}$ is $\le O(1)$ while the action
on the complement is $\ge\Omega(2^{4j}/(j+1)^{2})\gg 1$).

Applying the log-concave tail bound (Lemma~\ref{lem:log_concave_tail})
with $K=U_{\mathrm{AF},j}^{\mathrm{corr}}$ and $t=r_{\mathrm{AF},j}$:
\begin{equation*}
  \nu_{j}\bigl(\mathcal{Q}_{j}\setminus U_{\mathrm{AF},j}^{\mathrm{
  corr}}\bigr)
  \;\le\;2\,\exp\!\bigl(-c_{\mathrm{LC}}\,\kappa_{\mathrm{YM}}\,
  r_{\mathrm{AF},j}^{2}\bigr)
  \;=\;2\,\exp\!\bigl(-c_{\mathrm{LC}}\,\kappa_{\mathrm{YM}}\,
  (\eta_{\mathrm{AF}}^{\mathrm{corr}})^{2}\,2^{4j}/(j+1)^{2}\bigr).
\end{equation*}
Using $\kappa_{\mathrm{YM}}=c_{\mathrm{horiz}}/(2g_{\mathrm{YM}}^{2})$
and $g_{\mathrm{YM}}^{-2}(j)\sim\log(2^{j})\sim j\log 2$:
\begin{equation*}
  \nu_{j}\bigl(\mathcal{Q}_{j}\setminus U_{\mathrm{AF},j}^{\mathrm{
  corr}}\bigr)
  \;\le\;2\,\exp\!\bigl(-c\,(\eta_{\mathrm{AF}}^{\mathrm{corr}})^{2}\,
  2^{4j}\,j\log 2/(j+1)^{2}\bigr)
  \;\le\;c_{\mathrm{AF}}\,\exp\!\bigl(-\alpha_{\mathrm{AF}}\,
  2^{2j}/(j+1)\bigr),
\end{equation*}
where the last step uses $2^{4j}\,j/(j+1)^{2}\gg 2^{2j}/(j+1)$ for
$j\ge 1$ (and the $j=0$ case is absorbed into $c_{\mathrm{AF}}$).
\end{proof}

\subsection{The mass gap}
\label{ssec:app231_mass_gap}

We now combine the gauge-invariant concentration
(Theorem~\ref{thm:concentration_main}) with the corrected bracket
bound, Hessian coercivity, BE curvature, and per-scale Cheeger bound
of app230 (which all hold on $U_{\mathrm{AF},j}^{\mathrm{corr}}$) to
obtain the mass gap.

\begin{theorem}[Mass gap]
\label{thm:mass_gap_app231}
The Yang--Mills theory on $\mathbb{T}^{4}$ with gauge group
$G=\mathrm{SU}(N)$, $N\ge 2$, has a mass gap: there exists
$m_{\mathrm{phys}}>0$ such that the spectrum of the transfer matrix
has a gap $1-\lambda_{1}\ge c_{0}\,a_{n}^{2}$ (in lattice units) for
all sufficiently large $n$, giving
$m_{\mathrm{phys}}=(1-\lambda_{1})/a_{n}\ge c_{0}\,a_{n}\to 0$.
\end{theorem}

\begin{proof}[Proof sketch]
By the gauge-invariant concentration
(Theorem~\ref{thm:concentration_main}), the measure $\nu_{j}$
concentrates on $U_{\mathrm{AF},j}^{\mathrm{corr}}$ with exponential
rate $\alpha_{\mathrm{AF}}\,2^{2j}/(j+1)$. On
$U_{\mathrm{AF},j}^{\mathrm{corr}}$, the corrected bracket bound
(Lemma~\ref{lem:bracket_bound_corr}), Hessian coercivity
(Theorem~\ref{thm:hessian_coercivity_corr}), BE curvature bound
(Theorem~\ref{thm:be_curvature_corr}), and per-scale Cheeger bound
(Theorem~\ref{thm:sharp_cheeger_corr}) all hold, giving
$h_{j}^{\nu}\ge c_{h}\,2^{j}/(j+1)$ on
$U_{\mathrm{AF},j}^{\mathrm{corr}}$.

The multiscale aggregation (with the corrected normalization from
app230, accounting for the $a_{n}$-dependence of the OS transfer-matrix
Cheeger constant) gives the OS Cheeger constant
$h_{\mathrm{OS}}^{(n)}\ge c_{\mathrm{OS}}\,a_{n}^{-1/2}$ (in lattice
units), corresponding to $h_{\mathrm{OS}}^{(n)}\sim c_{\mathrm{OS}}\,
a_{n}^{1/2}$ in physical units. By the Cheeger bound for the transfer
matrix (Theorem~\ref{thm:cheeger_bound} of app218), the spectral gap
$1-\lambda_{1}^{(n)}\ge(c_{\mathrm{OS}})^{2}\,a_{n}/(2\mathfrak{D}_{
\mathrm{phys}}\,L_{\mathrm{phys}}^{2})$, giving the physical mass gap
$m_{\mathrm{phys}}=(1-\lambda_{1}^{(n)})/a_{n}\ge(c_{\mathrm{OS}})^{
2}/(2\mathfrak{D}_{\mathrm{phys}}\,L_{\mathrm{phys}}^{2})>0$.

The concentration estimate \eqref{eq:concentration_main} enters the
multiscale aggregation as follows: the contribution of the complement
$\mathcal{Q}_{j}\setminus U_{\mathrm{AF},j}^{\mathrm{corr}}$ to the
Cheeger constant is bounded by the concentration
$\nu_{j}(\mathcal{Q}_{j}\setminus U_{\mathrm{AF},j}^{\mathrm{corr}})
\le\exp\!\bigl(-\alpha_{\mathrm{AF}}\,2^{2j}/(j+1)\bigr)$, which is
exponentially small in $j$ and therefore does not affect the
polynomial lower bound $h_{j}^{\nu}\ge c_{h}\,2^{j}/(j+1)$ (the
exponential loss is absorbed into the constant $c_{h}$).
\end{proof}

\begin{remark}[Status of the Clay Millennium Prize problem]
\label{rem:clay_status}
The mass gap (Theorem~\ref{thm:mass_gap_app231}) establishes the
existence of $\Delta_{\mathrm{YM}}>0$ such that
$\inf\mathrm{spec}(-\Delta_{\mathrm{YM}})\ge\Delta_{\mathrm{YM}}>0$.
This is the mass-gap portion of the Clay Millennium Prize problem
on Yang--Mills theory. The remaining Clay requirements (existence of
the quantum field theory, conformal invariance at short distances,
and the full axiomatic framework) are addressed in the companion
appendices (app190--app228) and are not affected by the
concentration argument of this appendix.
\end{remark}

\begin{remark}[Key new ingredients]
\label{rem:new_ingredients}
The proof of the gauge-invariant concentration
(Theorem~\ref{thm:concentration_main}) introduces two new ingredients
not present in app229--app230:
\begin{enumerate}
\item \textbf{The Faddeev--Popov Gaussian lower bound on $Z_{j}$}
(Lemma~\ref{lem:gaussian_lower_Z}): a non-circular lower bound on
$Z_{j}$ obtained by integrating in a fixed gauge (Landau gauge) near
the flat connection, where the Hessian is the gauge-fixed Laplacian.
This breaks the circular dependency of Remark~\ref{rem:circular}.
\item \textbf{The log-concavity of $\nu_{j}$ on $U_{\mathrm{AF},j}^{
\mathrm{corr}}$} (Theorem~\ref{thm:log_concave}): a consequence of
the Hessian coercivity (Theorem~\ref{thm:hessian_coercivity_corr}),
which allows the log-concave concentration inequality
(Lemma~\ref{lem:log_concave_tail}) to be applied, giving the
exponential tail bound directly.
\end{enumerate}
The combination of these two ingredients closes the gauge-invariant
concentration hypothesis (Hypothesis~\ref{hyp:gauge_invariant_concentration})
and completes the mass-gap proof.
\end{remark}

\endinput
